%%%%%%%%%%  scan idx 227  |  folio 214  %%%%%%%%%%
Identifiant $u$, par (3.1.4), à une application continue de $S$ dans $\Diff^n(E,F)$, on constate que, pour $U$ ouvert dans $S$, le morphisme
\[
  \pr_{1*}(u)(U):\Gamma(X|U,E)\longrightarrow \Gamma(X|U,F)
\]
s'identifie canoniquement à l'application linéaire
\[
  C(U,C^\infty(V,E))\longrightarrow C(U,C^\infty(V,F))
\]
déduite de
\begin{equation}\tag{3.1.3}
  \varphi u:U\longrightarrow \Lop(C^\infty(V,E),C^\infty(V,F)).
\end{equation}
Supposons $S$ localement compact. Alors (3.1.1) le morphisme
\[
  \pr_{1*}(u):\pr_{1*}\pr_2^*(E)
  \longrightarrow \pr_{1*}\pr_2^*(F)
\]
est le morphisme de faisceaux associé au morphisme de préfaisceaux
\[
  u_\infty:\underline C'(S,C^\infty(V,E))
  \longrightarrow \underline C'(S,C^\infty(V,F))
\]
défini par $\varphi u$. D'autre part, pour $r\in\Z$,
\[
  \varphi_r u:S\longrightarrow \Lop(C^r_2(V,E),C^{r-n}_2(V,F))
\]
définit un morphisme de préfaisceaux
\[
  u_r:\underline C'(S,C^r_2(V,E))
  \longrightarrow \underline C'(S,C^r_2(V,F)).
\]
D'après (3.1.3 (ii)), les $u_r$ forment un système projectif dont la limite est $u_\infty$. Il nous sera commode d'interpréter $\varphi_r u$ comme un morphisme de fibrés hilbertiens triviaux (3.0.5)
\[
  S\times C^r_2(V,E)\longrightarrow S\times C^{r-n}_2(V,F),
\]
